\section{Lifting over small extensions}\label{sec:lifting}

The following obstruction calculation applies to isolated canonical
singularities beyond the local models of Theorem~\ref{thm:main}.

\begin{lemma}\label{lem:obstruction}
Let $X$ be a normal projective complex threefold with isolated Gorenstein
canonical singularities, $\omega_X\simeq\OO_X$ and $H^1(X,\OO_X)=0$.
Let $A'\twoheadrightarrow A$ be a small extension of local Artin
$\C$-algebras and let $X_A$ be a flat deformation of $X$.
If the restriction of $X_A$ at each singular point admits a lift to $A'$,
then $X_A$ admits a lift to $A'$.
\end{lemma}

\begin{proof}
Write $I$ for the kernel of the extension, so $\mathfrak m_{A'}I=0$.
We first explain the use of projective embeddings on the given Artin
deformation. Serre duality gives $H^2(X,\OO_X)=0$. The exact sequence of
units for a square-zero thickening therefore lifts any line bundle on $X$
to $X_A$, by induction on the length of $A$. For a very ample $M$ with
$H^1(X,M)=0$, a basis of $H^0(X,M)$ lifts to sections generating
the lifted bundle. They give a morphism $f_A:X_A\to\PP^N_A$ inducing
the embedding $f$ by $M$. This is a closed immersion: it is proper
with finite fibres, hence finite, and the map on structure sheaves is
surjective by Nakayama's lemma. In particular an analytic Artin
deformation can be treated algebraically, by GAGA for its ideal in
$\PP^N_A$. Local analytic lifts induce lifts of the completed local
rings. In comparing obstruction classes there we complete a
finite-type cotangent presentation, using continuous differentials.
Completion commutes with the finite free terms through each fixed
degree and is faithfully flat on the local obstruction module.

Put $L_X=L_{X/\C}$. The canonical obstruction to lifting $X_A$ is
\[
 o\in\Ext^2_X(L_X,\OO_X)\otimes I.
\]
Indeed the relative obstruction with coefficients in $I\otimes_\C\OO_X$ has
this target by flat base change for the cotangent complex and derived
adjunction. It is the composite of the transitivity boundary for
$X_A\to\Spec A\to\Spec A'$ with the shifted projection of
$L_{A/A'}$ onto $I[1]$. Its vanishing is equivalent to a flat lift
\cite[Tags 08UZ and 063Y]{Stacks}.
For the morphism $f_A$ above, functoriality of the transitivity triangle
identifies the image of $o$ under
\[
 \Ext^2_X(L_X,\OO_X)\otimes I
    \longrightarrow H^2(X,f^*T_{\PP^N})\otimes I
\]
with the pullback of the obstruction of $\PP^N_A$. Thus this image is
zero.

The negative cohomology sheaves of $L_X$ are supported on the finite
singular set $Z$. Finite-length modules over the local rings at these
points have grade three. The truncation of $L_X$ to $\Omega_X^1$ hence
induces an isomorphism in $\Ext^2(-,\OO_X)$, globally and locally
\cite[Lemma~2.2]{W5}. By Serre duality, $o$ is an $I$-valued functional
on $H^1(X,\Omega_X^1)$. The last vanishing says that this functional
kills the image of
\[
 H^1(X,f^*\Omega^1_{\PP^N})\longrightarrow H^1(X,\Omega_X^1).
\]
The connecting class of $1$ in the pulled-back Euler sequence maps to
$d\log M$, up to sign. Every line bundle is the difference, in
$\operatorname{Pic}(X)$, of two very ample bundles with the required
vanishing. Consequently $o$ annihilates $d\log\operatorname{Pic}(X)$.

By local duality, localization on $\Ext^2$ is dual to the map
\[
 H^1_Z(X,\Omega_X^1)\longrightarrow H^1(X,\Omega_X^1)
\]
\cite[Lemma~2.4(c)]{Gross}. At a singular point the image of $o$
with coefficients in the completed local ring is obtained by
precomposing the obstruction of the completed germ with the map
of cotangent complexes. The given local lift makes it zero, by the
same functoriality. Thus $o$ annihilates this second image. Write $j:X\setminus Z\hookrightarrow X$ and
$F=j_*\Omega^1_{X\setminus Z}$. The quotient by the displayed image is
$H^1(X,F)$. Indeed $H^1(X,F)$ embeds in $H^1(X\setminus Z,\Omega^1)$
by Leray, and $H^1(X,\Omega_X^1)\to H^1(X,F)$ is surjective since the
kernel and cokernel on sheaves have finite support. The local cohomology
sequence gives the assertion. Gross proves that the images of
$d\log\operatorname{Pic}(X)$ span $H^1(X,F)$
\cite[Lemma~2.5]{Gross}. Thus $o=0$.
\end{proof}

\begin{proof}[Proof of Theorem~\ref{thm:main}]
Write $B=\Def(X)$, let $G\subset V$ be the product of the marked local
period images, and put $T=\C^h\times(K\times_VG)$, with the indicated
scheme structure. Proposition~\ref{prop:period} gives the analytic
period map $P:B\to K\times_VG$. Its differential maps onto $K$ with
kernel $H^1(X,\Theta_X)$. Choose $h$ analytic functions $z$ vanishing
at the origin whose differentials restrict to a basis of the dual of this kernel. Then $F=(z,P):B\to T$ is a closed immersion of germs. Indeed
$dF$ is an isomorphism at the origin, and functions spanning the
cotangent space give a closed immersion by the analytic inverse
function theorem in a minimal smooth embedding.

Let $A'\twoheadrightarrow A$ be a small extension with kernel $I$.
Suppose $b\in B(A)$ and $y'\in T(A')$ restricts to $F(b)$. In the
construction of Proposition~\ref{prop:period}, lift the auxiliary
ambient parameter of $b$ to $A'$, which is possible since its ambient
germ is smooth. Over that lift, the inverse local period isomorphisms,
applied to the period coordinates of $y'$, give local hull parameters
extending those of $b$. In particular every singular germ of the
deformation over $b$ lifts to $A'$. Lemma~\ref{lem:obstruction} gives
a global lift, and the formal versality of $B$ gives an $A'$-point
$b'$ above $b$.

The difference of two lifts of the same point through a small
extension, viewed on local rings, is a derivation at the closed point
with values in $I$. Thus $y'-F(b')$ belongs to $T_0T\otimes I$.
Translate $b'$ by its inverse image under $dF\otimes 1_I$. The
resulting lift maps to $y'$. This proves formal smoothness of the
closed immersion $F$. More explicitly, successive lifting on the
Artin quotients of $\OO_{T,0}$ gives a section of $F$ on completions.
The completed quotient homomorphism is therefore an isomorphism.
Faithful flatness of completion gives the asserted analytic
isomorphism. The identifications on local branches and their marked
persistence loci follow from Proposition~\ref{prop:period}.
\end{proof}
